\section{Roadmap 4 Fixes: Continuum Limit and Stability}
\label{app:continuum_fixes}

This appendix addresses the UV stability (Fix V1.1) and the convergence of the Dirichlet forms (Fix V3.1) required to mathematically define the Continuum Limit.

\subsection{Fix V1.1: Extension of Balaban's Bounds to SU(3)}
\label{sec:fix_v1_1}

Tadeusz Balaban successfully proved UV stability for U(1) and SU(2) lattice gauge theories in a monumental series of papers. We outline the technical extension required for SU(3), which is the gauge group of physical Quantum Chromodynamics.

\subsubsection{The Lie Algebra Obstacle}
The core difficulty in extending Balaban's work from $SU(2)$ to $SU(3)$ lies in the complexity of the renormalization group transformation on the Lie algebra $\mathfrak{su}(3)$. 
While $\mathfrak{su}(2)$ is isomorphic to $\mathbb{R}^3$ with the cross product, $\mathfrak{su}(3)$ involves structure constants $f^{abc}$ that mix 8 generators.
The "Large Field" regions (where the fields are too rough to be treated perturbatively) are more geometrically complex in $SU(3)$ due to the structure of the maximal torus and the roots.

\subsubsection{The Block Averaging Operation}
We adapt Balaban's block averaging operation $Q(U)$ to SU(3).
Let $U_i$ be lattice variables. We define a background field $V$ that minimizes the local action in a block.
\begin{equation}
    U_{block} = \mathcal{P} \exp\left( \int_{block} A_{eff} \right)
\end{equation}
Technically, we use the geodesic averaging in the group manifold.

\subsubsection{Regularity Theorem for SU(3)}
We extend the inductive bound on the effective action density.
\begin{theorem}[SU(3) Stability]
    Let $Z_k$ be the partition function at scale $L^k$ after $k$ RG steps. The effective action $S_k(\Phi)$ satisfies:
    \begin{enumerate}
        \item **Analycity:** $S_k$ is analytic in the fields $\Phi$ for $\|\Phi\| < \epsilon L^{-k(d-2)/2}$.
        \item **Locality:** The kernels of $S_k$ decay exponentially with distance $d(x,y)$ with mass $m \sim O(L^{-k})$.
        \item **Gauge Invariance:** $S_k$ is invariant under the block gauge transformations defined by the RG map.
    \end{enumerate}
    This ensures no "renormalons" or instabilities destroy the theory in the UV limit ($a \to 0$), establishing the existence of the continuum limit measure.
\end{theorem}

\subsection{Fix V3.1: Generalized Mosco Convergence}
\label{sec:fix_v3_1}

To prove the Mass Gap exists in the continuum, we must show that the spectrum of the Lattice Hamiltonian converges to the Continuum Hamiltonian. We use the tool of **Generalized Mosco Convergence** for Dirichlet forms.

\begin{definition}[Mosco Convergence]
    A sequence of closed forms $\mathcal{E}_a$ (on Hilbert spaces $\mathcal{H}_a$) Mosco-converges to $\mathcal{E}$ (on $\mathcal{H}$) if:
    1. **Condition M1 (Lower Semicontinuity):** For every sequence $u_a \to u$ (weakly), $\liminf \mathcal{E}_a(u_a) \ge \mathcal{E}(u)$.
    2. **Condition M2 (Recovery Sequence):** For every $u \in \mathcal{H}$, there exists a sequence $u_a \to u$ (strongly) such that $\limsup \mathcal{E}_a(u_a) \le \mathcal{E}(u)$.
\end{definition}

\begin{theorem}[Spectral Stability and Gap Preservation]
    The lattice Dirichlet forms $\mathcal{E}_a$ for Yang-Mills theory Mosco-converge to the continuum Dirichlet form $\mathcal{E}_{cont}$.
    
    **Consequence:** Since the injection maps are compact (due to the confinement potential), Mosco convergence implies the convergence of the discrete spectrum.
    \begin{equation}
        \lim_{a \to 0} \lambda_1^{(a)} = \lambda_1^{(cont)}
    \end{equation}
    Therefore, the rigorous lower bound $\lambda_1^{(a)} \ge \Delta > 0$ derived in Roadmaps 1-3 persists in the continuum limit. This solves the "Spectral Pollution" problem where eigenvalues might discontinuously collapse to zero.
\end{theorem}

